\chapter{Dataset}
\label{Chapter3}
\lhead{Chapter 3. \emph{Dataset Description and Clinical Scenarios}}

\section{Overview}

The AgentClinic framework converts static medical datasets into Objective Structured Clinical Examination (OSCE) scenarios, enabling LLM evaluation in a dynamic, interactive setting. While the broader framework conceptually supports extensions to real-world EHR databases, this thesis strictly bounds its experimental scope to USMLE-style vignettes to maintain standardized evaluation.

For all experiments, a curated subset of the \textbf{AgentClinic-MedQA} dataset was utilized, comprising \textbf{107 high-quality diagnostic scenarios}. This subset was chosen for its breadth of clinical presentation, diversity of organ systems, and a representative distribution of diagnostic difficulty.

\section{The OSCE Examination Structure}

Unlike a static prompt, each scenario is parsed into a structured JSON format that enforces strict knowledge partitioning across agent roles. The Doctor agent must extract information through multi-turn dialogue and targeted test requests, rather than reasoning from a fully-specified vignette. The OSCE schema contains the following modules.

\subsection{Objective for Doctor}

A string describing the consultation objective (e.g., \textit{``Evaluate and diagnose the patient presenting with chest pain and shortness of breath.''}). This is the \textit{only} information the Doctor agent receives at the start of each case.

\subsection{Patient Actor Sub-schema}

This data is provided exclusively to the Patient Agent prompt to ground its dialogue:
\begin{itemize}
    \item \textbf{Demographics:} Patient age, gender, and relevant social background.
    \item \textbf{History of Present Illness (HPI):} A narrative detailing symptom onset, progression, and associated factors.
    \item \textbf{Symptoms:} Primary and secondary symptom descriptors (e.g., \textit{pain improves upon sitting forward; no productive cough}).
    \item \textbf{Past Medical and Social History:} Chronic conditions, prior surgeries, and lifestyle habits including smoking and alcohol use.
\end{itemize}

\subsection{Measurement Findings Sub-schema}

This information is held exclusively by the Measurement Agent and is disclosed only when the Doctor makes an explicit, well-formed test request:
\begin{itemize}
    \item \textbf{Vital Signs:} Temperature, blood pressure, heart rate, and respiratory rate.
    \item \textbf{Physical Examination:} Inspection, auscultation, and palpation findings.
    \item \textbf{Investigation Results:} Electrocardiogram, chest radiograph, blood panels (CBC, BMP), and imaging summaries.
\end{itemize}

\subsection{Ground Truth}

The verified \textbf{Correct Diagnosis} is provided strictly to the Moderator Agent, which evaluates the Doctor's final output by semantic equivalence comparison. This design prevents ground-truth contamination of any active agent role.

\section{Semantic Evaluation by the Moderator Agent}

Rather than enforcing exact string matching, which would penalise clinically valid terminological variation, the Moderator Agent judges diagnostic correctness by evaluating the semantic equivalence between the Doctor's stated diagnosis and the ground-truth label. It outputs a binary response (``Yes'' / ``No'') after comparing the prediction to the reference. This LLM-mediated evaluation accommodates synonymy, partial phrasing, and clinically acceptable variant naming conventions~\cite{schmidgall2024agentclinic}.
